\chapter{ASS1ST round-1 input (menu 22)}
\label{ch:menu22}
\menulabel{22}

\section{Purpose}

The first round of the ASS1ST active-space construction (Khedkar and Roemelt,
JCTC 2019/2020): a CASSCF calculation whose perturbation-theory density is
kept, so that menu~\menu{23} can read the quasi-natural occupation numbers and
propose the next active space.  The scheme builds an active space bottom-up
from a small but chemically reasonable starting point, which is what this menu
generates.

\section{What you need}

A structure in XYZ format and the initial active space $(n_{\mathrm{e}},
n_{\mathrm{o}})$ -- small but chemically reasonable, in the source's sense:
singly occupied orbitals for open shells, the metal $d$ shell for first-row
transition metals, a few $\pi/\pi^{*}$ orbitals for conjugated systems.  The
outcome depends on this choice, but with conservative thresholds different
sensible initials converge together (the source's own observation).

\section{Walkthrough}

\lstinputlisting[style=fbkcode, firstline=54, lastline=80,
  title={The menu-22 example session (water at an initial (4,4) space: the
  generated input carries the FIC-NEVPT2 + KeepDens keywords and the
  \code{PTSettings} block the round needs)}]
  {examples/expected/m22_ass1st_start.txt}

\section{What you get}

\file{<stem>.r1.inp} (the round-1 input) and \file{<stem>.r1.json.conf} (the
export request), both written next to the structure file, plus the run and
export guidance printed on screen.

The generated input carries the blocks the round needs, each measured on ORCA
6.1.1: \code{FIC-NEVPT2} (the unrelaxed density exists only for the FIC
ansatz; SC is refused), \code{KeepDens} (the density is read from the run's
\code{<base>.densities} sidecar), and the CASSCF \code{PTSettings} block with
\code{Density Unrelaxed} and \code{NatOrbs true}.

\section{How to read it}

\begin{readbox}
\begin{itemize}
  \item run the input (\code{orca <stem>.r1.inp}), export it
    (\code{orca\_2json <stem>.r1.gbw}; the \code{.json.conf} sits beside the
    \code{.gbw}), and feed \file{<stem>.r1.json} to menu~\menu{23};
  \item \textbf{the round's active orbitals are selected by ORCA's own default
    window from its starting guess} -- the source's near-ideal quasi-natural
    restart (feeding the previous round's rotated orbitals back) is not
    written yet.  Check the active orbitals in menu~\menu{23}'s output (it
    prints their CASSCF occupations): if the window missed the intended set,
    adjust the input before continuing the chain;
  \item the number of states decides the round's character: more than one
    makes it state-averaged, and menu~\menu{23} then averages the per-state
    densities with the state weights (equal by default);
  \item subsequent rounds are written by menu~\menu{23} as
    \file{<stem>.r2}, \file{<stem>.r3},~\dots\ -- the chain keeps the
    structure's name.
\end{itemize}
\end{readbox}
